\documentclass[portuguese,12pt,a4paper]{article}
\usepackage[T1]{fontenc}
\usepackage{graphicx}
\usepackage{mathtools}
\usepackage{amssymb}
\usepackage{amsthm}
\usepackage{thmtools}
\usepackage{xcolor}
\usepackage{nameref}
\usepackage{babel}
\usepackage{hyperref}
\title{Análise dos dados abertos da Agência Nacional do Cinema - ANCINE}
\author{Joamerson Islan}
\begin{document}
	\maketitle
	
	\section{Terminologia}
	\subsection{Certificado de Produto Brasileiro (CPB)}
	O Certificado de Produto Brasileiro (CPB) é o documento que regulariza as obras audiovisuais não publicitárias brasileiras.
	\subsection{Tipos de Obra}
	\subsubsection{Comum}
	O tipo comum engloba as obras audiovisuais que não estão contidas no espaço total do canal de programação, chamado de "espaço qualificado". Exclui-se do espaço qualificado:
	\begin{itemize}
		\item conteúdos religiosos ou políticos;
		\item manifestações e eventos esportivos;
		\item concursos;
		\item publicidade;
		\item televendas;
		\item infomerciais;
		\item jogos eletrônicos; propaganda política obrigatória;
		\item conteúdo audiovisual veiculado em horário eleitoral gratuito;
		\item conteúdos jornalísticos e programas de auditório ancorados
		por apresentador
	\end{itemize}
	Portanto, classifica-se como "comum" qualquer obra cinematográfica que seja de um dos tipos supracitados.
	\subsubsection{Brasileira constituinte de espaço qualificado}
	Esse tipo engloba produções que estão contidas no escopo do espaço qualificado, mas cujos direitos patrimoniais sobre a obra \textbf{são de posse de empresa ligada ao ramo audiovisual}.
	\subsubsection{Brasileira independente constituinte de espaço qualificado}
	Esse tipo engloba produções que estão contidas no escopo do espaço qualificado, cujos direitos patrimoniais sobre a obra \textbf{não são de posse de empresa ligada ao ramo audiovisual}.
	
	\subsection{Organização temporal das obras}
	\subsubsection{Obras seriadas}
	Obras audiovisuais seriadas são aquelas que são produzidas em capítulos ou episódios, ainda que com o mesmo título. Esse tipo recebe ainda algumas subclassificações, a saber:
	\begin{itemize}
		\item Temporada única;
		\item Múltiplas temporadas;
		\item Duração indeterminada.
	\end{itemize}
	\subsubsection{Obras não seriadas}
	Obras audiovisuais não seriadas são aquelas que simplesmente não se enquadram nas obras seriadas.
	
	\subsection{Segmento de mercado}
	O segmento de mercado é uma fração do espaço econômico audiovisual através da qual obras podem ser veiculadas a fim de levá-las a um determinado público-alvo. Os segmentos reconhecidos pela ANCINE são:
	\begin{itemize}
		\item Salas de Exibição;
		\item Radiodifusão de Sons e Imagens (TV Aberta);
		\item Comunicação Eletrônica de Massa por Assinatura (TV Paga);
		\item Vídeo Doméstico;
		\item Vídeo por Demanda;
		\item Outros;
		\subitem Audiovisual em Circuito Restrito;
		\subitem Audiovisual em Transporte Coletivo.
	\end{itemize}
	
	
	
\end{document}